\section{限界と倫理}
\label{sec:limitations-ethics}

\emph{分析の範囲．}
すべての分析は凍結したエンコーダと線形ヘッドを用い，微調整の確認は，低ランクのアダプタを用いた1つの音声エンコーダを4つのターゲットと3回の反復で調べたものである．
2つのガウス分布のモデルは，各感情の対の中でスコア差が共通の分散をもつガウス分布に従うと仮定し，式\eqref{eq:lambda}，式\eqref{eq:value}，命題\ref{prop:design}は1次または2次の近似に依拠する．
本論文は厳密な結果を主張せず，それらの予言を実験で検証する．
判断基準が先に直るという説明が成り立つのは単一モダリティの音声と顔の特徴で$K \le 5$の場合であり，$K=5$の融合の構成とテキストの構成では成り立たず，2つの損失がほぼ等しい顔で判断基準の割合が大きい理由はまだ説明できていない．

\emph{データ．}
英語の6つのコーパスから作った7つの群の4つの感情を調べ，コーパスごとにラベルの定義が異なる（\ref{sec:setting-data}節）．
CMU-MOSEIでは話者の分離が不完全であり，RAVDESSには本論文の処理の流れで顔の特徴が無い．
テキストを含む構成には4つの群しかないので，テキストを含む手法の比較は4つのターゲットにもとづき，第2のエンコーダは，提案法についてはAVTの構成でだけ検証した．

\emph{統計と選択．}
本論文の区間と検定は，本論文のコーパスの中でのラベルとテスト話者の引き方による変動を表す（\ref{sec:setting-selection}節）．
仮のターゲットでハイパーパラメータを選ぶと，$K=0$ではテスト話者での最良の値に対して最大で\NUM{cmp-oursx-oracle-vs-tuned-avt-t-k0}ポイントを失う（\ref{sec:comparison-selection}節）．
列の正規化は，本論文のターゲットのプールで感情の数が偏っているにもかかわらず（表\ref{tab:data}．たとえばMELD），感情ごとの合計をそろえる．
列の正規化による利得がクラスの比率にどう依存するかは分析しておらず，偏りがさらに強いと列の正規化は害になりうる\cite{veilleux2021realistic}．

\emph{倫理．}
公開されたコーパスだけをその配布条件の下で使い，新たなデータは記録していない．
顔は，検出した顔の切り出しの埋め込みとしてだけ処理の流れに入り，個人を識別できるデータは公開しない．
感情認識は，監視や，本人の同意のない人についての決定に悪用されうる．
本論文の手法は演技とメディアの録画からカテゴリのラベルを推定するものであり，個人についての決定のためには検証されていない．
